\section{Introduction}

Attention contains two matrix products separated by softmax.  NVIDIA
Blackwell executes 4-bit floating-point (FP4) matrix multiplication much
faster than bfloat16 (BF16), but FP4 does not accelerate the work between
those products.  Softmax must still reduce each score tile, evaluate
exponentials, construct a scaled probability tile, and make that tile ready
for the second product.  This middle stage becomes the bottleneck as the
matrix products get faster.

Attention combines a query matrix $Q$, key matrix $K$, and value matrix $V$:
\begin{equation}
  S=QK^\mathsf{T}/\sqrt d,\qquad
  P=\softmax(S),\qquad
  O=PV.
  \label{eq:attention-intro}
\end{equation}
We call the two matrix products QK and PV.  Here $S$ contains similarity
scores, $P$ contains the resulting probabilities, and $O$ is the output.
FlashAttention evaluates these equations in tiles so that the quadratic score
and probability matrices need not be stored in high-bandwidth memory (HBM)
\cite{dao2022flashattention}.  Later versions improved how work is divided and
overlapped \cite{dao2023flashattention2,shah2024flashattention3}.
FlashAttention-4 (FA4) adapts this tiled algorithm to NVIDIA Blackwell's
asynchronous matrix hardware \cite{zadouri2026flashattention4}.

We study forward inference and causal training separately.  In the forward
pass, we ask whether all four operands inside attention---$Q$, $K$, $P$, and
$V$---can use FP4 and still outperform BF16 FA4.  We build on HAO AI Lab's
two-query Blackwell schedule and replace its probability path
\cite{hao2026fp4flashattention,zhang2026attnqat}.  Our method, \emph{Direct-P},
maps normalized scores directly to MXFP4 probability codes and normalizes the
same rounded values that the PV product consumes.  In this paper, ``full FP4''
refers only to these four attention operands; it does not mean that every
transformer operation uses FP4.

In training, we ask whether causal backward can reuse the low-precision state
created by the forward pass.  Forward passes the quantized $Q/K$ payload,
its scales, and the softmax normalizer directly to backward.  The projection
and gradient epilogues also publish the row- or column-oriented FP8 views
needed by each gradient matrix product.  Backward reconstructs the
probabilities from this state instead of creating a separate BF16 score path.
Learned projections remain a separate precision boundary, so this is not
pure-FP4 end-to-end training.

The paper makes three contributions:
\begin{enumerate}
  \item Direct-P reaches up to 2.13$\times$ the BF16 forward throughput on
  favorable Blackwell shapes; matched fixed-input evaluations measure its
  error against an FP8 probability path.
  \item Passing the forward quantization into backward accelerates
  projection-inclusive attention and a
  single-GPU 8-billion-parameter update.  Distributed diagnostics select FP8
  P/V because every tested MXFP4 P/V trajectory diverges.
  \item Hardware profiles identify tensor-memory ownership as the main limit on
  overlap and motivate another allocatable score destination with compatible
  issue semantics.
\end{enumerate}

The public implementation and reproduction materials are linked in
Appendix~\ref{app:reproduction}.

Table~\ref{tab:reader-notation} defines the notation and hardware terms used
throughout the paper.  E$x$M$y$ names a floating-point format with $x$
exponent bits and $y$ explicit fraction bits; for example, E4M3 is an 8-bit
format and E2M1 is the payload used by the FP4 formats studied here.

\begin{table}[htbp]
\centering
\caption{Reader's guide to the main notation and Blackwell hardware terms.}
\label{tab:reader-notation}
\small
\begin{tabularx}{\textwidth}{p{.21\textwidth}X}
\toprule
Term & Meaning \\
\midrule
$B,S,H,H_q,H_{kv},D$ & Batch size, sequence length, head count, query-head
count, key/value-head count, and per-head dimension. \\
Shape shorthand & Compact labels append each value: B1/S4096/H24/D128 means
batch 1, sequence length 4096, 24 heads, and head dimension 128. \\
FP32, BF16, FP8, FP4 & 32-, 16-, 8-, and 4-bit floating-point families.
Smaller formats increase matrix throughput but need explicit scaling. \\
NVFP4 and MXFP4 & The two block-scaled FP4 families used here.  NVFP4 uses
fine-grained data-dependent scales; MXFP4 shares one power-of-two scale across
each 32-value block. \\
SM and CTA & A graphics processing unit (GPU) contains streaming
multiprocessors (SMs).  A cooperative thread array (CTA) is a CUDA thread
block scheduled on an SM. \\
TMEM & Tensor memory: Blackwell's on-chip accumulator scratchpad for
asynchronous matrix operations. \\
TMA and MMA & The Tensor Memory Accelerator (TMA) moves tiles; a matrix
multiply--accumulate (MMA) instruction performs the tensor-core product. \\
\bottomrule
\end{tabularx}
\end{table}

Sections~\ref{sec:hao-background}--\ref{sec:direct-p} develop the forward
argument in order: inherited schedule, remaining problems, and proposed
fixes.  Section~\ref{sec:causal-training} keeps isolated backward,
projection-inclusive attention, complete model updates, and distributed
training trajectories as separate measurement boundaries.
Section~\ref{sec:current-headroom} then relates the measured bottlenecks to
hardware and closes with the evidence boundaries and conclusions.
